\chapter{102}

Valentina verliess das Gebäude.

Da hatte sie wohl gerade etwas Bedeutungsvolles erlebt. Ganz von selbst
schüttelte sich ihr Körper. Sie hielt sich am Treppengeländer fest.

Dann zog sie die warme Luft tief in sich ein und schaute sich um.

Der nächstliegende, schattige Ort lag nur einige Meter von ihr entfernt.
Keiner ging direkt daran vorbei. Sie brauchte Ruhe.

Sie wählte die direkte Nummer des Commissario.

Eine Frauenstimme meldete sich. Valentina war kurz irritiert.

«Einen Moment, bitte, Signora, Commissario Mantovani ist gerade am
Telefon,» hörte sie sie sagen, «Sie können dran bleiben.»

«Si, grazie.»

Valentinas Herz klopfte bis zum Hals. Sie liess sich in einem Kreis von
Sitzsteinen nieder. Es dauerte.

Gerade als sie den Anruf beenden wollte, hörte sie ihn.

«Signora Ricci, sind Sie noch da?»

Mantovanis warme Stimme tat ihr gut.

«Si, Commissario, Buongiorno.»

«Buongiorno, Signora.»

«Ich war bei Giorgia Sala im \emph{Ospedale}, ich bin gerade
rausgekommen und dachte mir, ehm, also sie ist wach.»

«Ja, Signora, seit gestern. Das ist wunderbar. Ein Arzt des «\emph{San
Raffaele»} rief mich vorhin an. Sie sind die Erste, mit der sie
offensichtlich bewusst gesprochen hat.»

«Ja, ein Krankenpfleger hat mir erklärt, dass Giorgia seit gestern
spreche, aber bis vorhin nur wenige Worte. Ich, ich freue mich, es war
schön.»

«Sie scheinen einen guten Einfluss auf sie zu haben, Signora.»

«Danke, Commissario. Ja, wir konnten tatsächlich miteinander reden. Sie
sprach von Bastiano di Martino, also, sie glaube, er sei vor dem Absturz
bei ihr gewesen.»

«Wie bitte? Was meinen Sie mit… \emph{sie glaube}?»

«Giorgia ist sich offenbar nicht sicher.»

«Si.»

«Und, ja, ich fragte, ob denn sonst noch jemand da gewesen sei. Und sie
verneinte, vehement, wie mir schien. Irgend etwas murmelte sie noch von
einem Ast.»

«Signora Ricci, es ist sehr gut, dass Sie sich melden. Ich danke Ihnen.
Haben Sie die Möglichkeit, im Verlauf des Nachmittags auf die Questur zu
kommen? Ich möchte Ihre Aussage gerne offiziell aufnehmen.»

«Ja, Commissario, ich komme bei Ihnen vorbei.»

Sie verabschiedeten sich. Und Valentina wusste nicht genau, warum sie
sich wie angeschossen durch und durch elend fühlte.
